% GENERATED FILE. Edit canonical lessons, then run npm run generate:books.
% canonical-source-hash: fnv1a64:0e337a45f788c15e
% canonical-lessons: BN-C533-doibat, BN-C533-kaktaliyo, BN-C533-mooka, BN-C533-phoskano, BN-C533-lotari

\chapter{By Chance}
\label{ch:bn-by-chance}
\hlchaptermodality{533}

\begin{chapteropening}
\textbf{By the end of this chapter:} \emph{I can talk about chance meetings, coincidences and missed chances.}

Complete the last lesson of chapter 533: I can talk about chance meetings, coincidences and missed chances.
\end{chapteropening}

\section[\texorpdfstring{\bn{দৈবাৎ} (doibāt)}{দৈবাৎ (doibāt)}]{\bn{দৈবাৎ} (doibāt) — by chance}

\label{lesson:BN-C533-doibat}



\begin{quote}
Before the new one: say the Bengali for a forecast, then the Bengali for an astrologer.
\end{quote}

\subsection*{You'll want to know: \bn{দৈবাৎ}}
\textbf{\bn{দৈবাৎ}} — \emph{doibāt} — \textquotedblleft{}by chance\textquotedblright{}.

\begin{quote}
\bn{বাজারে} \bn{দৈবাৎ} \bn{আমার} \bn{পুরনো} \bn{বন্ধুর} \bn{সঙ্গে} \bn{দেখা} \bn{হল।}
\end{quote}

\begin{quote}
\emph{Bājāre doibāt āmār purôno bôndhur sôngge dækhā hôlô.}
\end{quote}

\begin{quote}
By chance I met my old friend at the market.
\end{quote}

\begin{cousinweb}
\textbf{\bn{দৈবাৎ}} (\emph{doibāt}) is the Sanskrit form \emph{daivāt}, \textquotedblleft{}from fate\textquotedblright{}, of \emph{daiva}, \textquotedblleft{}coming from the gods, fate\textquotedblright{} (Monier-Williams): by the gods' doing.
\end{cousinweb}

\begin{grammarlens}[title={What you've built}]
One more word for saying what happens if.
\end{grammarlens}

\subsection*{Your turn}
\begin{itemize}
\raggedright
  \item \emph{Say it:} \emph{doibāt}
  \item \emph{Say it:} \emph{doibāt}, once more
  \item \emph{Say it:} \emph{Bājāre doibāt āmār purôno bôndhur sôngge dækhā hôlô.} — By chance I met my old friend at the market.
  \item \emph{Recall:} say the Bengali for a forecast, then the Bengali for an astrologer, then say \emph{doibāt} again
\end{itemize}

\begin{tcolorbox}[breakable,colback=teal!4,colframe=teal!35!black,title={Before you move on}]
What does \bn{দৈবাৎ} mean? (\textquotedblleft{}By chance\textquotedblright{}.) Say it once more.
\end{tcolorbox}
\section[\texorpdfstring{\bn{কাকতালীয়} (kāktālīyô)}{কাকতালীয় (kāktālīyô)}]{\bn{কাকতালীয়} (kāktālīyô) — coincidental}

\label{lesson:BN-C533-kaktaliyo}



\begin{quote}
Before the new one: say the Bengali for an astrologer, then the Bengali for by chance.
\end{quote}

\subsection*{You'll want to know: \bn{কাকতালীয়}}
\textbf{\bn{কাকতালীয়}} — \emph{kāktālīyô} — \textquotedblleft{}coincidental\textquotedblright{}.

\begin{quote}
\bn{আমাদের} \bn{জন্মদিন} \bn{একই} \bn{দিনে}, \bn{এটা} \bn{কাকতালীয়।}
\end{quote}

\begin{quote}
\emph{Āmāder jônmôdin eki dine, eṭā kāktālīyô.}
\end{quote}

\begin{quote}
Our birthdays fall on the same day; it's a coincidence.
\end{quote}

\begin{cousinweb}
\textbf{\bn{কাকতালীয়}} (\emph{kāktālīyô}) is Sanskrit \emph{kākatālīya} (Monier-Williams), from a fable: a crow lands on a palm tree at the very moment a fruit falls, and the two events only look connected.
\end{cousinweb}

\begin{grammarlens}[title={What you've built}]
One more word for saying what happens if.
\end{grammarlens}

\subsection*{Your turn}
\begin{itemize}
\raggedright
  \item \emph{Say it:} \emph{kāktālīyô}
  \item \emph{Say it:} \emph{kāktālīyô}, once more
  \item \emph{Say it:} \emph{Āmāder jônmôdin eki dine, eṭā kāktālīyô.} — Our birthdays fall on the same day; it's a coincidence.
  \item \emph{Recall:} say the Bengali for an astrologer, then the Bengali for by chance, then say \emph{kāktālīyô} again
\end{itemize}

\begin{tcolorbox}[breakable,colback=teal!4,colframe=teal!35!black,title={Before you move on}]
What does \bn{কাকতালীয়} mean? (\textquotedblleft{}Coincidental\textquotedblright{}.) Say it once more.
\end{tcolorbox}
\section[\texorpdfstring{\bn{মওকা} (môokā)}{মওকা (môokā)}]{\bn{মওকা} (môokā) — a lucky break; an opening}

\label{lesson:BN-C533-mooka}



\begin{quote}
Before the new one: say the Bengali for by chance, then the Bengali for coincidental.
\end{quote}

\subsection*{You'll want to know: \bn{মওকা}}
\textbf{\bn{মওকা}} — \emph{môokā} — \textquotedblleft{}a lucky break; an opening\textquotedblright{}.

\begin{quote}
\bn{দোকানে} \bn{ছাড়} \bn{ছিল}, \bn{আমি} \bn{মওকা} \bn{বুঝে} \bn{একটা} \bn{জ্যাকেট} \bn{কিনলাম।}
\end{quote}

\begin{quote}
\emph{Dokāne chhāṛ chhilô, āmi môokā bujhe ekṭā jækeṭ kinlām.}
\end{quote}

\begin{quote}
There was a discount at the shop; I saw my chance and bought a jacket.
\end{quote}

\begin{cousinweb}
\textbf{\bn{মওকা}} (\emph{môokā}) is usually traced to Arabic \emph{mawqiʿ}, \textquotedblleft{}a place, an occasion\textquotedblright{}, through Persian.
\end{cousinweb}

\begin{grammarlens}[title={What you've built}]
One more word for saying what happens if.
\end{grammarlens}

\subsection*{Your turn}
\begin{itemize}
\raggedright
  \item \emph{Say it:} \emph{môokā}
  \item \emph{Say it:} \emph{môokā}, once more
  \item \emph{Say it:} \emph{Dokāne chhāṛ chhilô, āmi môokā bujhe ekṭā jækeṭ kinlām.} — There was a discount at the shop; I saw my chance and bought a jacket.
  \item \emph{Recall:} say the Bengali for by chance, then the Bengali for coincidental, then say \emph{môokā} again
\end{itemize}

\begin{tcolorbox}[breakable,colback=teal!4,colframe=teal!35!black,title={Before you move on}]
What does \bn{মওকা} mean? (\textquotedblleft{}A lucky break; an opening\textquotedblright{}.) Say it once more.
\end{tcolorbox}
\section[\texorpdfstring{\bn{ফসকানো} (phôskāno)}{ফসকানো (phôskāno)}]{\bn{ফসকানো} (phôskāno) — to slip away; to be missed (a chance)}

\label{lesson:BN-C533-phoskano}



\begin{quote}
Before the new one: say the Bengali for coincidental, then the Bengali for a lucky break.
\end{quote}

\subsection*{You'll want to know: \bn{ফসকানো}}
\textbf{\bn{ফসকানো}} — \emph{phôskāno} — \textquotedblleft{}to slip away; to be missed (a chance)\textquotedblright{}.

\begin{quote}
\bn{এই} \bn{সুযোগ} \bn{ফসকালে} \bn{তুমি} \bn{আর} \bn{পাবে} \bn{না।}
\end{quote}

\begin{quote}
\emph{Ei sujog phôskāle tumi ār pābe nā.}
\end{quote}

\begin{quote}
If you let this chance slip, you won't get it again.
\end{quote}

Past (he/she …): \bn{সেই} \bn{সুযোগ} \bn{ফসকাল}, \emph{sei sujog phôskālô}.

\begin{cousinweb}
\textbf{\bn{ফসকানো}} (\emph{phôskāno}) probably belongs with \bn{ফসকা}, \textquotedblleft{}loose\textquotedblright{}, which CDIAL lists as Bengali \emph{phaskā} under a reconstructed \emph{phass}, \textquotedblleft{}loose\textquotedblright{} (CDIAL 9068): what slips from a loose grip, a ball or a chance.
\end{cousinweb}

\begin{grammarlens}[title={What you've built}]
One more word for saying what happens if.
\end{grammarlens}

\subsection*{Your turn}
\begin{itemize}
\raggedright
  \item \emph{Say it:} \emph{phôskāno}
  \item \emph{Say it:} \emph{phôskāno}, once more
  \item \emph{Say it:} \emph{Ei sujog phôskāle tumi ār pābe nā.} — If you let this chance slip, you won't get it again.
  \item \emph{Recall:} say the Bengali for coincidental, then the Bengali for a lucky break, then say \emph{phôskāno} again
\end{itemize}

\begin{tcolorbox}[breakable,colback=teal!4,colframe=teal!35!black,title={Before you move on}]
What does \bn{ফসকানো} mean? (\textquotedblleft{}To slip away; to be missed (a chance)\textquotedblright{}.) Say it once more.
\end{tcolorbox}
\section[\texorpdfstring{\bn{লটারি} (lôṭāri)}{লটারি (lôṭāri)}]{\bn{লটারি} (lôṭāri) — a lottery}

\label{lesson:BN-C533-lotari}



\begin{quote}
Before the new one: say the Bengali for a lucky break, then the Bengali for to slip away.
\end{quote}

\subsection*{You'll want to know: \bn{লটারি}}
\textbf{\bn{লটারি}} — \emph{lôṭāri} — \textquotedblleft{}a lottery\textquotedblright{}.

\begin{quote}
\bn{লটারিতে} \bn{জিতলে} \bn{মামা} \bn{একটা} \bn{নতুন} \bn{গাড়ি} \bn{কিনবেন।}
\end{quote}

\begin{quote}
\emph{Lôṭārite jitle māmā ekṭā nôtun gāṛi kinben.}
\end{quote}

\begin{quote}
If my uncle wins the lottery, he will buy a new car.
\end{quote}

\begin{cousinweb}
\textbf{\bn{লটারি}} (\emph{lôṭāri}) is borrowed from English \emph{lottery}.
\end{cousinweb}

\begin{grammarlens}[title={What you've built}]
One more word for saying what happens if.
\end{grammarlens}

\subsection*{Your turn}
\begin{itemize}
\raggedright
  \item \emph{Say it:} \emph{lôṭāri}
  \item \emph{Say it:} \emph{lôṭāri}, once more
  \item \emph{Say it:} \emph{Lôṭārite jitle māmā ekṭā nôtun gāṛi kinben.} — If my uncle wins the lottery, he will buy a new car.
  \item \emph{Recall:} say the Bengali for a lucky break, then the Bengali for to slip away, then say \emph{lôṭāri} again
\end{itemize}

\begin{tcolorbox}[breakable,colback=teal!4,colframe=teal!35!black,title={Before you move on}]
What does \bn{লটারি} mean? (\textquotedblleft{}A lottery\textquotedblright{}.) Say it once more.
\end{tcolorbox}
